\documentclass[11pt]{article}
\usepackage[a4paper,margin=25mm]{geometry}
\usepackage{amsmath,amssymb,amsthm}
\usepackage[T1]{fontenc}
\usepackage{lmodern}
\usepackage[hidelinks]{hyperref}
\newtheorem{theorem}{Theorem}
\newcommand{\Q}{\mathbb Q}
\newcommand{\Z}{\mathbb Z}
\newcommand{\Gal}{\operatorname{Gal}}
\title{A Fixed Integer Counterexample Family\newline for Conjecture 00000000107}
\author{GitHub: gaochengzhecpu}
\date{}
\begin{document}
\maketitle
\begin{abstract}
The fixed nonzero integer $a=2$ disproves the proposed eventual $S_n$ Galois group for $x^n-x-a$. For every even $n\ge2$, the rational root $-1$ implies that the actual Galois group over $\Q$ has order at most $(n-1)!<n!$. We prove an unbounded family of counterexamples, including the impossibility of an abstract group isomorphism with $S_n$.
\end{abstract}

\section*{Statement and relation to Conjecture 95}
The source says that for every fixed nonzero integer $d$, as $n\to\infty$, the Galois group of $x^n-x-d$ over $\Q$ is $S_n$. Writing $a$ for the fixed integer parameter, its usual eventual interpretation is
\[
 \forall a\in\Z\setminus\{0\}\ \exists N\ \forall n\ge N,\quad
 \Gal(x^n-x-a\,/\Q)\cong S_n.
\]
Both source languages include $a=2$ and impose no parity restriction on $n$.

The same fixed-constant rational-root obstruction also applies to the prime-parameter statement in Conjecture 00000000095. We explicitly reuse that construction and its general Galois-group lemmas here. This submission addresses the integer-parameter quantifier of Conjecture 107 and contains its own complete proof and Lean project; it has no dependency on another submission or PR.

\begin{theorem}\label{thm:counter}
For every even integer $n\ge2$, let $f_n(x)=x^n-x-2\in\Q[x]$. Then
\[
 \bigl|\Gal(f_n/\Q)\bigr|\le(n-1)!<n!.
\]
Thus $\Gal(f_n/\Q)$ is not isomorphic to $S_n$.
\end{theorem}
\begin{proof}
Evenness gives $f_n(-1)=(-1)^n+1-2=0$. Therefore $f_n=(x+1)q_n$ for a rational polynomial $q_n$ of degree $n-1$. Its splitting field already contains $-1$, so it is also a splitting field of $f_n$. The two polynomials consequently have isomorphic Galois groups.

Let $R$ be the set of distinct roots of $q_n$ in its splitting field. Every automorphism over $\Q$ permutes $R$. This action is faithful, because the splitting field is generated by those roots over $\Q$. Hence
\[
 \bigl|\Gal(f_n/\Q)\bigr|
 =\bigl|\Gal(q_n/\Q)\bigr|
 \le |R|!\le(n-1)!.
\]
Since $n\ge2$, one has $(n-1)!<n!=|S_n|$. Isomorphic finite groups have equal orders, proving the assertion.
\end{proof}

For every threshold $N\ge0$, the even degree $n=2(N+1)$ is at least $N$ and at least $2$. It therefore contradicts the asserted conclusion at the fixed nonzero integer $a=2$. No eventual threshold can exist for that parameter. The argument proves failure at every positive even degree, not merely at a few small degrees. It also excludes validity on a set of degrees of density one. We do not claim anything here about odd degrees or a modified statement that excludes the parameter $2$.

\section*{Formal verification}
The independent project \texttt{lean/Main.lean} uses Mathlib's actual rational polynomial type and \texttt{Polynomial.Gal}, the group of rational algebra automorphisms of the polynomial's splitting field. The constant parameter of \texttt{trinomial} has type $\Z$, and its degree parameter has type $\mathbb N$.

The formal proof verifies the complete chain:
\begin{itemize}
\item The polynomial is $X^n-X-C(a)$, with proved degree $n$ for $n\ge2$. For $a=2$, its value at $-1$ vanishes at every even degree.
\item The factor theorem produces an actual rational polynomial quotient. Its degree is $n-1$, obtained from the degree formula for a product of nonzero polynomials.
\item Mathlib's faithful root action bounds the order of any polynomial's Galois group by the factorial of its degree. This requires no irreducibility assumption on the quotient.
\item The injective restriction map for the Galois group of a product, together with triviality of the group of a rational linear factor, gives the required bound for $f_n$. This supplies the formal inequality directly, without a separate formal identification of splitting fields.
\item Group orders rule out a multiplicative equivalence with \texttt{Equiv.Perm (Fin n)}. This excludes an abstract group isomorphism; the conclusion does not rely on a particular root labeling.
\item The final result is \texttt{conjecture107\_false}. It negates the eventual assertion quantified over \emph{all nonzero integers}, by taking the parameter $2$ and constructing a counterexample beyond every threshold.
\end{itemize}
The eight printed theorem audits concern these generic results. The optional Python script checks the root, explicit factorization, quotient degree, and factorial comparison for six sample even degrees. It does not compute exact Galois groups or use finite experiments to infer the infinite conclusion. The Lean conclusion is the pointwise even-degree family and the negation of the eventual assertion; no separate density formalization is claimed.

\section*{Source and reproduction}
The original bilingual source is preserved byte-for-byte in \texttt{SOURCE.md}. Its repository path is
\begin{center}\small\href{https://github.com/The-Last-Math-Competition/The-Last-Math-Competition/blob/main/conjectures/00000000107.md}{\texttt{conjectures/00000000107.md}}.\end{center}
The project pins Lean 4.19.0 and Mathlib commit
\begin{center}\small\texttt{c44e0c8ee63ca166450922a373c7409c5d26b00b}.\end{center}
From \texttt{lean/}, run \texttt{lake build} followed by
\begin{center}\small\texttt{lake env lean -DwarningAsError=true Main.lean}.\end{center}
Dependency URLs are public Git URLs with locked revisions. Actual clean-build logs, theorem-axiom audits, source hashes, and PDF checks are included under \texttt{verification/}; \texttt{README.md} supplies the remaining reproduction instructions.
\end{document}
